\section{Proof Debugging: Turning Fog Into Work}
\chapterline

A proof rarely arrives clean on the first try. Even experienced mathematicians write scratch work, notice a missing condition, rename variables, split a case, or replace a vague sentence by a definition. The difference between a beginner and a mature proof writer is not that the mature writer never gets stuck. The mature writer has a method for finding the stuck point.

This section is about that method.

\begin{focusbox}
When a proof feels wrong, do not ask first, ``What is the trick?'' Ask, ``What is the first sentence whose justification I cannot defend?''
\end{focusbox}

Debugging a proof is a disciplined way to convert uneasiness into work. You identify the claim, unpack the grammar, locate the first unsupported step, and repair that step with a definition, a lemma, a witness, a case split, or a counterexample.

\subsection{The proof feedback workflow}

Use this workflow whenever you read your own proof, review someone else's proof, or get feedback from an instructor.

\begin{enumerate}
\item \textbf{Restate the claim.} Write the theorem in one sentence, with all domains visible.
\item \textbf{Mark the quantifiers.} Circle every ``for all,'' ``there exists,'' ``for some,'' ``unique,'' and ``there is no.''
\item \textbf{Separate hypotheses from conclusion.} Put assumptions on one side and the target on the other.
\item \textbf{Type every object.} Record whether each symbol is an integer, real number, set, function, relation, subset, sequence, or something else.
\item \textbf{Open the target definition.} If the conclusion says ``even,'' write ``equals $2k$ for some $k\in\Z$.'' If it says ``injective,'' write the exact implication to prove.
\item \textbf{Open the useful hypothesis definitions.} Do not open every definition automatically; open the ones that can produce the shape needed by the target.
\item \textbf{Choose the outer proof form.} Direct proof, contrapositive, contradiction, cases, induction, or two inclusions.
\item \textbf{Write the skeleton.} Before details, write the first and last two lines.
\item \textbf{Fill one bridge at a time.} Every line should be licensed by algebra, a definition, a hypothesis, a lemma, or a previous line.
\item \textbf{Run the failure check.} Look for undefined variables, reversed implications, hidden division by zero, lost cases, and conclusions stronger than what was proved.
\end{enumerate}

\begin{formalbox}
A proof line is healthy when a reader can answer both questions:
\[
\text{What does this line claim?}
\qquad
\text{Why is this line allowed?}
\]
If either answer is unclear, the proof needs repair.
\end{formalbox}

\subsection{Proof Anatomy Worksheet}

Before debugging a proof, dissect it. This is the proof equivalent of symbolic execution.

\begin{longtable}{p{0.22\linewidth}p{0.36\linewidth}p{0.32\linewidth}}
\toprule
Anatomy item & Question & Example prompt \\
\midrule
Objects & What mathematical things are involved? & integers, sets, functions, groups \\
Domains & Where do the objects live? & $n\in\Z$, $f:A\to B$ \\
Assumptions & What may be used? & $f$ is injective, $A\subseteq B$ \\
Target & What must be shown? & $g\circ f$ is injective \\
Definitions & Which words must be opened? & injective, subset, compact \\
Tools & Which previous facts apply? & contrapositive, induction, De Morgan \\
Strategy & What proof shape fits? & direct, cases, contradiction \\
Intermediate states & What bridge claims are needed? & $f(a_1)=f(a_2)$ before $a_1=a_2$ \\
Closure & Where is the target satisfied? & ``therefore by definition...'' \\
\bottomrule
\end{longtable}

\begin{center}
\begin{tikzpicture}[
  node distance=0.75cm,
  box/.style={draw, rounded corners, align=center, fill=gray!5, minimum width=4.2cm, minimum height=0.7cm},
  arrow/.style={-{Latex}, thick}
]
\node[box] (a) {Assumptions};
\node[box, below=of a] (d) {Definitions opened};
\node[box, below=of d] (i) {Intermediate states};
\node[box, below=of i] (t) {Target closed};
\draw[arrow] (a) -- (d);
\draw[arrow] (d) -- (i);
\draw[arrow] (i) -- (t);
\node[right=1.0cm of d, text width=4.4cm] {If a proof feels foggy, locate which arrow is unsupported.};
\end{tikzpicture}
\end{center}

\begin{labbox}
Use the worksheet on this claim: if $f:A\to B$ is injective and $S,T\subseteq A$, then
\[
f(S\cap T)=f(S)\cap f(T).
\]
Identify where injectivity is needed. Then write a counterexample showing the claim can fail without injectivity.
\end{labbox}

\subsection{The first-bad-line rule}

When reviewing a proof, do not merely say that the proof is wrong. Find the first line where the argument stops being justified. Later lines may also be false, but the first bad line is the place where repair begins.

For example, consider the claim:

\begin{quote}
If $n^2$ is even, then $n$ is even.
\end{quote}

Bad proof:

\begin{quote}
Since $n^2$ is even, $n^2=2k$. Taking square roots gives $n=\sqrt{2k}$, so $n$ is even.
\end{quote}

The first two moves are not the real problem. If $n^2$ is even, then there is some $k\in\Z$ such that $n^2=2k$. Taking square roots is already delicate because $n$ might be negative, but even if we replace it by $|n|=\sqrt{2k}$, the final inference is invalid. The expression $\sqrt{2k}$ need not have the form $2m$ for an integer $m$.

\begin{trapbox}
Do not let familiar algebraic motion impersonate proof. A line can look like calculation and still fail to establish the definition you need.
\end{trapbox}

A repair uses the contrapositive:

\begin{proof}
We prove the contrapositive. Suppose $n$ is odd. Then $n=2m+1$ for some $m\in\Z$. Hence
\[
n^2=(2m+1)^2=4m^2+4m+1=2(2m^2+2m)+1.
\]
Since $2m^2+2m\in\Z$, $n^2$ is odd. Therefore, if $n^2$ is even, then $n$ is even.
\end{proof}

\subsection{A feedback vocabulary}

Good feedback is specific. Here are useful phrases.

\begin{longtable}{p{0.25\linewidth}p{0.31\linewidth}p{0.34\linewidth}}
\toprule
Feedback phrase & What it means & Repair question \\
\midrule
Undefined variable & A symbol appears before introduction & Where does this object live? \\
Quantifier mismatch & The proof proves the wrong order of choices & Who chooses first? \\
Missing witness & An existence claim has no constructed object & What object actually works? \\
Definition not closed & The proof gestures at a property but never matches its definition & What exact form must be shown? \\
Converse error & The proof uses $Q\Rightarrow P$ when only $P\Rightarrow Q$ is known & Which direction is available? \\
Case leak & One possibility is not covered & What are all exhaustive cases? \\
Circular step & The proof uses what it is trying to prove & Can this line be justified without the conclusion? \\
Type error & The expression combines objects of incompatible kinds & What kind of object is each symbol? \\
Hidden assumption & The proof uses a condition not in the theorem & Is this condition given or proved? \\
Too strong & The proof establishes or assumes something stronger than needed & What weaker target is enough? \\
\bottomrule
\end{longtable}

\subsection{How to comment on a bad proof}

Weak feedback:

\begin{quote}
This is confusing.
\end{quote}

Better feedback:

\begin{quote}
The variable $k$ is introduced in line 2, but its domain is not stated. If you mean $k\in\Z$, say so. The final line also needs to show that the expression has the form $2m$ for an integer $m$.
\end{quote}

Weak feedback:

\begin{quote}
You need more rigor.
\end{quote}

Better feedback:

\begin{quote}
The proof says $f$ is ``one-to-one by cancellation.'' To prove injectivity, start with arbitrary $a_1,a_2\in A$ and assume $f(a_1)=f(a_2)$. Then show $a_1=a_2$. The missing step is the cancellation calculation.
\end{quote}

\begin{intuitionbox}
Good proof feedback sounds like a bug report with a reproducible failure: here is the line, here is the expected definition, here is what is missing.
\end{intuitionbox}

\subsection{The repair loop}

Use the following loop when revising.

\begin{enumerate}
\item Copy the broken line.
\item Name the reason it is broken in one phrase.
\item Write the definition or theorem that should justify it.
\item Add the missing intermediate line.
\item Check whether the new line introduced a new obligation.
\item Repeat until every obligation is discharged.
\end{enumerate}

For example:

\begin{quote}
Broken line: ``Therefore $x\in A$ and $x\in B$.''\\
Problem: This was concluded from $x\in A\cup B$.\\
Definition: $x\in A\cup B$ means $x\in A$ or $x\in B$.\\
Repair: Use cases, or change the target.
\end{quote}

\subsection{Bad proof repair 1: parity}

Claim: If $a$ and $b$ are odd integers, then $ab$ is odd.

Bad proof:

\begin{quote}
Odd times odd is odd because an odd number has an odd factor, so the product is odd.
\end{quote}

Diagnosis:

\begin{itemize}
\item The proof repeats the claim instead of proving it.
\item The phrase ``has an odd factor'' does not close the definition of odd.
\item No variables witnessing oddness are introduced.
\end{itemize}

Repair:

\begin{proof}
Let $a,b\in\Z$ be odd. Then there exist $m,n\in\Z$ such that
\[
a=2m+1,\qquad b=2n+1.
\]
Thus
\[
ab=(2m+1)(2n+1)=4mn+2m+2n+1=2(2mn+m+n)+1.
\]
Since $2mn+m+n\in\Z$, the product $ab$ is odd.
\end{proof}

Commentary: The repaired proof opens the definition of odd twice and closes the definition of odd once.

\subsection{Bad proof repair 2: divisibility}

Claim: If $a\mid b$ and $b\mid c$, then $a\mid c$.

Bad proof:

\begin{quote}
Since $a$ divides $b$ and $b$ divides $c$, $a$ is part of $b$ and $b$ is part of $c$, so $a$ is part of $c$.
\end{quote}

Diagnosis:

\begin{itemize}
\item ``Part of'' is not a divisibility definition.
\item The proof gives a picture, not a witness.
\item The conclusion $a\mid c$ requires an integer $q$ such that $c=aq$.
\end{itemize}

Repair:

\begin{proof}
Suppose $a\mid b$ and $b\mid c$. Then there exist integers $m,n$ such that
\[
b=am,\qquad c=bn.
\]
Substituting $b=am$ into $c=bn$ gives
\[
c=(am)n=a(mn).
\]
Since $mn\in\Z$, we have $a\mid c$.
\end{proof}

\subsection{Bad proof repair 3: rational numbers}

Claim: The sum of two rational numbers is rational.

Bad proof:

\begin{quote}
Let $x=a/b$ and $y=c/d$. Then $x+y=(a+c)/(b+d)$, so the sum is rational.
\end{quote}

Diagnosis:

\begin{itemize}
\item The addition formula is false.
\item The proof does not state that $b,d\neq0$.
\item It accidentally produces an expression that can be undefined if $b+d=0$.
\end{itemize}

Repair:

\begin{proof}
Let $x,y\in\Q$. Then there exist integers $a,b,c,d$ with $b\neq0$, $d\neq0$, and
\[
x=\frac ab,\qquad y=\frac cd.
\]
Therefore
\[
x+y=\frac ab+\frac cd=\frac{ad+bc}{bd}.
\]
Since $ad+bc\in\Z$ and $bd\in\Z$ with $bd\neq0$, the number $x+y$ is rational.
\end{proof}

\begin{trapbox}
In proofs about rational numbers, the nonzero denominator is not decoration. It is part of the definition.
\end{trapbox}

\subsection{Bad proof repair 4: set equality}

Claim:
\[
A\cap(B\cup C)=(A\cap B)\cup(A\cap C).
\]

Bad proof:

\begin{quote}
This is true by distribution because intersection distributes over union.
\end{quote}

Diagnosis:

\begin{itemize}
\item This may be acceptable only if the distributive law has already been proved.
\item If the theorem is the distributive law, the proof is circular.
\item A beginner should prove it by element chasing.
\end{itemize}

Repair:

\begin{proof}
Let $x$ be arbitrary. Then
\[
x\in A\cap(B\cup C)
\]
if and only if $x\in A$ and $x\in B\cup C$. This holds if and only if $x\in A$ and either $x\in B$ or $x\in C$. By distributivity of logical ``and'' over logical ``or,'' this is equivalent to
\[
(x\in A\text{ and }x\in B)\quad\text{or}\quad(x\in A\text{ and }x\in C).
\]
That is equivalent to
\[
x\in(A\cap B)\cup(A\cap C).
\]
Since $x$ was arbitrary, the sets are equal.
\end{proof}

\subsection{Bad proof repair 5: image and preimage}

Claim:
\[
f^{-1}(S\cap T)=f^{-1}(S)\cap f^{-1}(T).
\]

Bad proof:

\begin{quote}
This is obvious because inverse functions preserve intersections.
\end{quote}

Diagnosis:

\begin{itemize}
\item The notation $f^{-1}(S)$ does not require $f$ to have an inverse function.
\item The proof does not unpack the definition of inverse image.
\item The word ``obvious'' is hiding all the work.
\end{itemize}

Repair:

\begin{proof}
Let $x$ be arbitrary. Then
\[
x\in f^{-1}(S\cap T)
\]
if and only if $f(x)\in S\cap T$. This holds if and only if $f(x)\in S$ and $f(x)\in T$. That is equivalent to saying
\[
x\in f^{-1}(S)\quad\text{and}\quad x\in f^{-1}(T).
\]
Therefore $x\in f^{-1}(S)\cap f^{-1}(T)$. Since membership in the two sets is equivalent for arbitrary $x$, the sets are equal.
\end{proof}

\subsection{Bad proof repair 6: injectivity}

Claim: If $f:A\to B$ and $g:B\to C$ are injective, then $g\circ f:A\to C$ is injective.

Bad proof:

\begin{quote}
Both functions are one-to-one, so no values collide. Therefore the composition is one-to-one.
\end{quote}

Diagnosis:

\begin{itemize}
\item The idea is correct but not yet a proof.
\item The proof does not begin with equal outputs of the composition.
\item It does not use injectivity in the standard direction.
\end{itemize}

Repair:

\begin{proof}
Let $a_1,a_2\in A$ and suppose
\[
(g\circ f)(a_1)=(g\circ f)(a_2).
\]
Then
\[
g(f(a_1))=g(f(a_2)).
\]
Since $g$ is injective, it follows that
\[
f(a_1)=f(a_2).
\]
Since $f$ is injective, it follows that $a_1=a_2$. Therefore $g\circ f$ is injective.
\end{proof}

\subsection{Bad proof repair 7: surjectivity}

Claim: If $f:A\to B$ and $g:B\to C$ are surjective, then $g\circ f:A\to C$ is surjective.

Bad proof:

\begin{quote}
Let $a\in A$. Since $f$ is onto, $f(a)=b$, and since $g$ is onto, $g(b)=c$. Therefore every $c$ is hit.
\end{quote}

Diagnosis:

\begin{itemize}
\item To prove surjectivity onto $C$, start with arbitrary $c\in C$, not arbitrary $a\in A$.
\item The proof uses $b$ and $c$ without correct quantifier order.
\item Surjectivity gives preimages of chosen outputs, not arbitrary symbolic equations.
\end{itemize}

Repair:

\begin{proof}
Let $c\in C$ be arbitrary. Since $g$ is surjective, there exists $b\in B$ such that $g(b)=c$. Since $f$ is surjective, there exists $a\in A$ such that $f(a)=b$. Then
\[
(g\circ f)(a)=g(f(a))=g(b)=c.
\]
Thus every $c\in C$ is hit by $g\circ f$, so $g\circ f$ is surjective.
\end{proof}

\subsection{Bad proof repair 8: contrapositive}

Claim: If $n^2$ is odd, then $n$ is odd.

Bad proof:

\begin{quote}
Assume $n^2$ is odd. Then $n^2=2k+1$. So $n=\sqrt{2k+1}$, which is odd.
\end{quote}

Diagnosis:

\begin{itemize}
\item Taking square roots does not produce the definition of odd.
\item The expression $\sqrt{2k+1}$ need not display an integer in the form $2m+1$.
\item The contrapositive is cleaner.
\end{itemize}

Repair:

\begin{proof}
We prove the contrapositive. Suppose $n$ is not odd. Since every integer is either even or odd, $n$ is even. Then $n=2m$ for some $m\in\Z$. Hence
\[
n^2=(2m)^2=4m^2=2(2m^2),
\]
so $n^2$ is even. Therefore, if $n^2$ is odd, then $n$ is odd.
\end{proof}

Commentary: A direct proof is possible, but it typically smuggles in a parity fact. Contrapositive proof uses the definition of even immediately.

\subsection{Bad proof repair 9: contradiction}

Claim: There are infinitely many prime numbers.

Bad proof:

\begin{quote}
Assume there are finitely many primes $p_1,\ldots,p_n$. Then $p_1p_2\cdots p_n+1$ is prime, so there is another prime. Contradiction.
\end{quote}

Diagnosis:

\begin{itemize}
\item The number $p_1p_2\cdots p_n+1$ need not be prime.
\item What is true is that no listed prime divides it.
\item The proof needs the fact that every integer greater than $1$ has a prime divisor.
\end{itemize}

Repair:

\begin{proof}
Suppose, for contradiction, that there are only finitely many primes, say
\[
p_1,p_2,\ldots,p_n.
\]
Let
\[
N=p_1p_2\cdots p_n+1.
\]
Since $N>1$, it has a prime divisor $q$. By the assumption that our list contains all primes, $q=p_i$ for some $i$. Then $q$ divides the product $p_1p_2\cdots p_n$, and $q$ also divides $N$. Hence $q$ divides their difference:
\[
N-p_1p_2\cdots p_n=1.
\]
This is impossible, since no prime divides $1$. Therefore there are infinitely many primes.
\end{proof}

\subsection{Bad proof repair 10: induction}

Claim:
\[
1+2+\cdots+n=\frac{n(n+1)}2
\]
for all $n\in\N$.

Bad proof:

\begin{quote}
It works for $1$. Assume it works for $n$. Then
\[
1+2+\cdots+n=\frac{n(n+1)}2.
\]
Therefore it works for $n+1$.
\end{quote}

Diagnosis:

\begin{itemize}
\item The induction step never proves the statement for $n+1$.
\item The proof repeats the induction hypothesis and stops.
\item The missing move is adding $n+1$ to both sides and simplifying to the correct target.
\end{itemize}

Repair:

\begin{proof}
For $n=1$, the left side is $1$ and the right side is $1(2)/2=1$.

Assume for some $k\in\N$ that
\[
1+2+\cdots+k=\frac{k(k+1)}2.
\]
Then
\[
1+2+\cdots+k+(k+1)=\frac{k(k+1)}2+(k+1).
\]
Factoring gives
\[
\frac{k(k+1)}2+(k+1)=\frac{k(k+1)+2(k+1)}2=\frac{(k+1)(k+2)}2.
\]
This is the desired formula with $k+1$ in place of $k$. Therefore the formula holds for all $n\in\N$.
\end{proof}

\subsection{Bad proof repair 11: uniqueness}

Claim: If an integer has an additive inverse, it is unique.

Bad proof:

\begin{quote}
The inverse of $a$ is $-a$, and there is only one $-a$.
\end{quote}

Diagnosis:

\begin{itemize}
\item This assumes the formula for the inverse rather than proving uniqueness from the inverse property.
\item A uniqueness proof should suppose there are two candidates and prove they are equal.
\item The relevant law is associativity and the identity property of $0$.
\end{itemize}

Repair:

\begin{proof}
Let $a\in\Z$. Suppose $b$ and $c$ are both additive inverses of $a$. Then
\[
a+b=0,\qquad a+c=0.
\]
Now
\[
b=b+0=b+(a+c)=(b+a)+c=(a+b)+c=0+c=c.
\]
Thus $b=c$, so the additive inverse is unique.
\end{proof}

Commentary: This proof pattern appears constantly in algebra: assume two objects have the universal property or defining property, then prove they must be equal.

\subsection{Bad proof repair 12: open sets}

Claim: The interval $(0,1)$ is open in $\R$.

Bad proof:

\begin{quote}
It is open because it does not contain its endpoints.
\end{quote}

Diagnosis:

\begin{itemize}
\item Endpoint language works for intervals in $\R$, but it is not the definition of open.
\item The proof must show that each point has a positive-radius neighborhood inside the interval.
\item The radius must depend on the point.
\end{itemize}

Repair:

\begin{proof}
Let $x\in(0,1)$. Define
\[
r=\frac12\min\{x,1-x\}.
\]
Since $0<x<1$, both $x$ and $1-x$ are positive, so $r>0$. If $y\in\R$ and $|y-x|<r$, then
\[
x-r<y<x+r.
\]
Because $r\leq x/2$, we have $x-r\geq x/2>0$. Because $r\leq(1-x)/2$, we have $x+r\leq x+(1-x)/2<1$. Hence $0<y<1$, so $y\in(0,1)$. Thus every point of $(0,1)$ has a ball contained in $(0,1)$, and $(0,1)$ is open.
\end{proof}

\subsection{Bad proof repair 13: limits}

Claim:
\[
\lim_{n\to\infty}\frac{4n+1}{n}=4.
\]

Bad proof:

\begin{quote}
As $n$ gets big, $1/n$ goes away, so the limit is $4$.
\end{quote}

Diagnosis:

\begin{itemize}
\item The intuition is correct, but the proof does not use the epsilon definition.
\item The quantity to control is
\[
\left|\frac{4n+1}{n}-4\right|.
\]
\item A final proof must choose $N$ in terms of $\eps$.
\end{itemize}

Repair:

\begin{proof}
Let $\eps>0$. Choose $N\in\N$ such that $N>1/\eps$. If $n\geq N$, then
\[
\left|\frac{4n+1}{n}-4\right|=\left|4+\frac1n-4\right|=\frac1n\leq\frac1N<\eps.
\]
Therefore $\frac{4n+1}{n}\to4$.
\end{proof}

\subsection{Bad proof repair 14: equivalence relations}

Claim: Congruence modulo $m$ is an equivalence relation on $\Z$, where $m\in\N$ and $m\geq2$.

Bad proof:

\begin{quote}
This is an equivalence relation because numbers with the same remainder are equivalent.
\end{quote}

Diagnosis:

\begin{itemize}
\item The sentence describes the idea but does not verify reflexive, symmetric, and transitive.
\item The proof should use the definition $a\equiv b\pmod m$ if $m\mid(a-b)$.
\item Each property needs its own argument.
\end{itemize}

Repair:

\begin{proof}
Define $a\sim b$ if $m\mid(a-b)$.

Reflexive: For any $a\in\Z$, we have $a-a=0=m\cdot0$, so $m\mid(a-a)$. Hence $a\sim a$.

Symmetric: Suppose $a\sim b$. Then $m\mid(a-b)$, so $a-b=mk$ for some $k\in\Z$. Thus
\[
b-a=-(a-b)=m(-k),
\]
so $m\mid(b-a)$, and $b\sim a$.

Transitive: Suppose $a\sim b$ and $b\sim c$. Then $a-b=mk$ and $b-c=m\ell$ for some $k,\ell\in\Z$. Adding gives
\[
a-c=(a-b)+(b-c)=m(k+\ell).
\]
Therefore $m\mid(a-c)$, so $a\sim c$.

Thus congruence modulo $m$ is an equivalence relation.
\end{proof}

\subsection{Bad proof repair 15: partial orders}

Claim: The subset relation $\subseteq$ is a partial order on $\pow(X)$.

Bad proof:

\begin{quote}
Subsets are ordered by size, so this is a partial order.
\end{quote}

Diagnosis:

\begin{itemize}
\item Subset inclusion is not the same as size comparison.
\item Two sets can have the same size without being equal.
\item A partial order requires reflexivity, antisymmetry, and transitivity.
\end{itemize}

Repair:

\begin{proof}
We check the three properties on $\pow(X)$.

Reflexive: For any $A\subseteq X$, every element of $A$ is an element of $A$, so $A\subseteq A$.

Antisymmetric: Suppose $A\subseteq B$ and $B\subseteq A$. Then every element of $A$ is in $B$, and every element of $B$ is in $A$. Hence $A=B$.

Transitive: Suppose $A\subseteq B$ and $B\subseteq C$. If $x\in A$, then $x\in B$, and therefore $x\in C$. Thus $A\subseteq C$.

Therefore $\subseteq$ is a partial order on $\pow(X)$.
\end{proof}

\subsection{Debugging drills}

For each proof below, identify the first bad line, name the gap, and write a repair.

\begin{enumerate}
\item Claim: If $x,y\in\Q$, then $xy\in\Q$. Bad proof: ``Since rationals are fractions, their product is a fraction.''
\item Claim: If $a\mid b$ and $a\mid c$, then $a\mid(b+c)$. Bad proof: ``Since $a$ goes into both, it goes into the sum.''
\item Claim: If $A\subseteq B$, then $A\cap C\subseteq B\cap C$. Bad proof: ``Intersections preserve inclusion.''
\item Claim: If $f$ is injective, then $f^{-1}(\{f(a)\})=\{a\}$. Bad proof: ``This follows because injective functions have inverses.''
\item Claim: If $n^2$ is divisible by $3$, then $n$ is divisible by $3$. Bad proof: ``Taking square roots preserves divisibility.''
\item Claim: If $a_n\to0$ and $b_n\to0$, then $a_n+b_n\to0$. Bad proof: ``Both are small eventually, so the sum is small.''
\item Claim: If $U$ and $V$ are open subsets of $\R$, then $U\cap V$ is open. Bad proof: ``The intersection of open things is open.''
\item Claim: If $R$ and $S$ are equivalence relations on $A$, then $R\cap S$ is an equivalence relation. Bad proof: ``Intersections preserve structure.''
\end{enumerate}

\begin{labbox}
For the next week, keep a proof error log. Each entry should have four lines: the broken sentence, the gap type, the repaired sentence, and the lesson. The goal is not to collect embarrassments. The goal is to make your mistakes searchable.
\end{labbox}

\section{The Velleman Readiness Plan}
\chapterline

Velleman's \emph{How to Prove It} is not a book of tricks. It is a book about proof architecture. You are ready to benefit from it when you can recognize the outer shape of a proof before knowing every inner move.

\begin{focusbox}
Readiness does not mean speed. It means you can locate the next responsible move.
\end{focusbox}

The plan below is deliberately detailed. It is meant to be used, not admired. If a week takes two weeks, that is fine. The point is to build a stable proof habit before the text accelerates.

\subsection{The readiness standard}

Before beginning Velleman seriously, you should be able to do most of the following without looking up a template.

\begin{enumerate}
\item Translate a quantified English sentence into symbols.
\item Translate a symbolic quantified sentence into clear English.
\item Negate statements with one, two, and three quantifiers.
\item Identify the hypothesis and conclusion of a conditional.
\item Prove a direct implication using definitions.
\item Decide when the contrapositive is cleaner than direct proof.
\item Use contradiction without leaving the contradiction vague.
\item Prove set equality by arbitrary element or two inclusions.
\item Prove subset claims by starting with an arbitrary element.
\item Prove injectivity and surjectivity from their definitions.
\item Work with images and inverse images of sets.
\item Verify that a relation is reflexive, symmetric, transitive, antisymmetric, or total.
\item Prove a relation is an equivalence relation.
\item Describe equivalence classes and quotient sets.
\item Write ordinary induction with a clean base case and induction step.
\item Recognize when strong induction is the natural tool.
\item Produce counterexamples to false universal claims.
\item Review your own proof using the first-bad-line rule.
\end{enumerate}

\subsection{A six-week readiness schedule}

The schedule assumes five study sessions per week. Each session should be about sixty to ninety minutes. If you have less time, preserve the order and reduce the number of exercises; do not skip the proof-revision days.

\begin{longtable}{p{0.08\linewidth}p{0.2\linewidth}p{0.34\linewidth}p{0.28\linewidth}}
\toprule
Week & Theme & Daily work & End-of-week standard \\
\midrule
1 & Mathematical sentences & Statements, predicates, domains, conditionals, negations & You can say exactly what a theorem claims \\
2 & Quantifiers & Translation, negation, order of choice, examples and counterexamples & You no longer treat $\forall\exists$ and $\exists\forall$ as interchangeable \\
3 & Core proof methods & Direct proof, contrapositive, contradiction, cases & You can choose a proof shape from the conclusion \\
4 & Sets and functions & Subsets, equality, power sets, image, inverse image, injective, surjective & You can run element-chasing proofs without guessing \\
5 & Relations and induction & Equivalence, partial order, classes, ordinary and strong induction & You can prove structural claims in organized parts \\
6 & Portfolio and exam & Rewrite, debug, timed proofs, oral explanations & You can begin Velleman with active reading habits \\
\bottomrule
\end{longtable}

\subsection{Daily session template}

Use this structure for most study days.

\begin{enumerate}
\item \textbf{Ten-minute recall.} Without notes, write definitions from yesterday.
\item \textbf{Twenty-minute reading.} Read slowly. Stop at each theorem and predict the proof shape.
\item \textbf{Thirty-minute proof writing.} Write one or two complete proofs in full sentences.
\item \textbf{Fifteen-minute debugging.} Review one proof using the feedback workflow.
\item \textbf{Five-minute log.} Record one definition, one mistake, and one question.
\end{enumerate}

\begin{walkbox}
At the end of a session, ask: what did I actually prove today, and which definition did I use most heavily?
\end{walkbox}

\subsection{Week 1: sentences and domains}

Goals:

\begin{itemize}
\item distinguish expressions, predicates, statements, commands, and questions;
\item identify the domain of every variable;
\item rewrite conditionals as ``if hypothesis, then conclusion'';
\item notice when a statement is true in one domain and false in another.
\end{itemize}

Core drills:

\begin{enumerate}
\item Decide whether each is a statement: $x^2=4$; ``there exists $x\in\R$ such that $x^2=4$''; ``solve $x^2=4$''; ``$n$ is prime.''
\item For $P(x):x^2<4$, determine the truth set over $\N$, $\Z$, $\Q$, and $\R$.
\item Rewrite: ``A square integer is nonnegative'' using a universal quantifier.
\item Rewrite: ``Some integer is larger than its square'' using an existential quantifier, then decide whether it is true.
\item Explain why the domain matters in the statement ``Every number has a multiplicative inverse.''
\end{enumerate}

Portfolio task:

\begin{quote}
Write a one-page note titled ``What a Mathematical Sentence Needs.'' Include five examples and three nonexamples.
\end{quote}

\subsection{Week 2: quantifier discipline}

Goals:

\begin{itemize}
\item negate quantified statements without guessing;
\item understand order of dependency;
\item distinguish pointwise and uniform claims;
\item build counterexamples to false universal statements.
\end{itemize}

Core drills:

\begin{enumerate}
\item Negate: $\forall x\in\R,\exists y\in\R,\ y>x$.
\item Negate: $\exists M\in\R,\forall x\in S,\ x\leq M$.
\item Negate: $\forall \eps>0,\exists N\in\N,\forall n\geq N,\ |a_n-L|<\eps$.
\item Explain the difference between ``every student has a favorite book'' and ``there is a book favored by every student.''
\item Give an example where $\forall x\exists y\,R(x,y)$ is true but $\exists y\forall x\,R(x,y)$ is false.
\end{enumerate}

Portfolio task:

\begin{quote}
Create a two-column page. In the left column write ten quantified statements. In the right column write their negations in English and symbols.
\end{quote}

\begin{trapbox}
The negation of ``for every $x$, there exists $y$'' is not ``for every $x$, there does not exist $y$.'' The correct negation is ``there exists $x$ such that for every $y$, not.''
\end{trapbox}

\subsection{Week 3: proof methods}

Goals:

\begin{itemize}
\item use direct proof when definitions flow forward;
\item use contrapositive when the negated conclusion has strong structure;
\item use contradiction when denying the conclusion creates a concrete object to defeat;
\item use cases only when the cases are exhaustive.
\end{itemize}

Core drills:

\begin{enumerate}
\item Prove directly: If $a,b$ are even integers, then $a+b$ is even.
\item Prove by contrapositive: If $n^2$ is even, then $n$ is even.
\item Prove by contradiction: $\sqrt2$ is irrational.
\item Prove by cases: For every integer $n$, $n^2+n$ is even.
\item For each proof above, write the proof skeleton before the final proof.
\end{enumerate}

Portfolio task:

\begin{quote}
Write four proofs of parity or divisibility claims: one direct, one contrapositive, one contradiction, one cases. Label every definition opened and closed.
\end{quote}

\subsection{Week 4: sets and functions}

Goals:

\begin{itemize}
\item prove subset claims by arbitrary elements;
\item prove set equality by membership equivalence or two inclusions;
\item distinguish image from inverse image;
\item prove injectivity and surjectivity using correct starting lines.
\end{itemize}

Core drills:

\begin{enumerate}
\item Prove $A\cap B\subseteq A$.
\item Prove $A\setminus B=A\cap B^c$.
\item Prove $f^{-1}(S\cup T)=f^{-1}(S)\cup f^{-1}(T)$.
\item Give a counterexample showing $f(S\cap T)=f(S)\cap f(T)$ can fail.
\item Prove that the composition of injective functions is injective.
\item Prove that the composition of surjective functions is surjective.
\end{enumerate}

Portfolio task:

\begin{quote}
Write a clean mini-chapter titled ``Element Chasing.'' Include one subset proof, one equality proof, one inverse-image proof, and one counterexample involving images.
\end{quote}

\subsection{Week 5: relations and induction}

Goals:

\begin{itemize}
\item check relation properties in separate paragraphs;
\item describe equivalence classes without handwaving;
\item understand quotient sets as sets of classes;
\item write induction proofs with an explicit induction statement.
\end{itemize}

Core drills:

\begin{enumerate}
\item Prove equality modulo $m$ is an equivalence relation on $\Z$.
\item Describe the equivalence classes modulo $4$.
\item Prove that subset inclusion is a partial order on $\pow(X)$.
\item Prove by induction that $1+3+\cdots+(2n-1)=n^2$.
\item Prove by induction that $2^n\geq n+1$ for all $n\geq1$.
\item Explain why strong induction is useful for proving every integer $n\geq2$ has a prime factorization.
\end{enumerate}

Portfolio task:

\begin{quote}
Write one full relation proof and one full induction proof. Then write a paragraph comparing the two proof structures.
\end{quote}

\subsection{Week 6: portfolio, review, and launch}

Goals:

\begin{itemize}
\item rewrite older proofs with cleaner definitions;
\item identify recurring personal errors;
\item take a timed readiness exam;
\item begin Velleman with an active reading protocol.
\end{itemize}

Core work:

\begin{enumerate}
\item Choose five old proofs and revise them.
\item For each revision, write the first bad line from the old version.
\item Complete the final exam at the end of this file under a time limit.
\item Explain three proofs out loud without looking.
\item Start Velleman Chapter 1 by translating every displayed symbolic statement into English.
\end{enumerate}

\subsection{The Velleman reading protocol}

When reading Velleman, do not read passively. Use the following procedure.

\begin{enumerate}
\item Before a section, skim headings and examples. Predict the skill being taught.
\item For each definition, write one example and one nonexample.
\item For each theorem, identify the proof shape before reading the proof.
\item After reading a proof, close the book and write the skeleton.
\item For each exercise, write the target definition before attempting the proof.
\item After solving, rewrite the proof once for clarity.
\item Keep a list of ``moves'' rather than isolated answers.
\end{enumerate}

\begin{historybox}
Velleman's book belongs to a tradition of proof texts that teach mathematics as language, strategy, and structure rather than as a list of finished theorems. Its exercises matter because they force you to practice the transitions: English to logic, logic to proof plan, proof plan to readable argument.
\end{historybox}

\subsection{Readiness self-test}

Answer these without notes.

\begin{enumerate}
\item What is the difference between a statement and a predicate?
\item What is the negation of $P\Rightarrow Q$?
\item How do you prove a universal statement?
\item How do you prove an existential statement?
\item How do you disprove a universal statement?
\item How do you prove $A\subseteq B$?
\item How do you prove $A=B$?
\item How do you prove a function is injective?
\item How do you prove a function is surjective?
\item What are the three properties of an equivalence relation?
\item What are the three properties of a partial order?
\item What are the three required parts of an induction proof?
\end{enumerate}

If you miss more than three, review before starting Velleman seriously. If you miss one or two, start but keep those topics in daily review.

\section{Second Pass and Mastery Review}
\chapterline

The first pass through proof teaches vocabulary. The second pass teaches control. You revisit the same ideas, but now you ask sharper questions:

\begin{itemize}
\item What is the exact logical form?
\item Which definition is doing the work?
\item What would a counterexample have to look like?
\item Which proof method fits the conclusion?
\item Where could a beginner reverse an implication?
\item How will this idea return in analysis, algebra, topology, or geometry?
\end{itemize}

\begin{focusbox}
Mastery is not remembering a theorem. Mastery is knowing what kind of work the theorem can do.
\end{focusbox}

\subsection{Mastery review method}

Use five passes for each major topic.

\begin{enumerate}
\item \textbf{Definition pass.} Write every definition from memory.
\item \textbf{Example pass.} Give three examples and two nonexamples.
\item \textbf{Proof pass.} Prove one standard theorem using the definition.
\item \textbf{Debug pass.} Repair one broken proof.
\item \textbf{Transfer pass.} Explain where the idea appears later.
\end{enumerate}

This method is slower than rereading. It is also much more honest. Rereading can feel fluent even when you cannot produce anything. Mastery review forces production.

\subsection{Proof Compression and Expansion}

A major barrier in advanced reading is proof compression. A mature text may write:

\begin{quote}
The result follows immediately from compactness.
\end{quote}

That phrase can hide an entire open-cover argument. Your job is to learn how to expand compressed moves until they become checkable.

There are several compression levels:

\begin{enumerate}
\item \textbf{Expanded proof:} every definition is opened and every witness is named.
\item \textbf{Textbook proof:} standard definitions are used fluently, but the strategy is visible.
\item \textbf{Proof sketch:} only the main mechanism is described.
\item \textbf{Expert shorthand:} standard lemmas and idioms are invoked by name.
\end{enumerate}

\begin{formalbox}
Expansion drill:
\begin{quote}
Whenever a proof says ``clearly,'' ``by definition,'' ``standard,'' or ``immediate,'' replace that phrase with three explicit lines.
\end{quote}
\end{formalbox}

Compressed:

\begin{quote}
Since $A\subseteq B$, clearly $\pow(A)\subseteq\pow(B)$.
\end{quote}

Expanded:

\begin{quote}
Let $S\in\pow(A)$. Then $S\subseteq A$. Since $A\subseteq B$, every element of $S$ is an element of $B$. Hence $S\subseteq B$, so $S\in\pow(B)$. Therefore $\pow(A)\subseteq\pow(B)$.
\end{quote}

\begin{trapbox}
False intuition: if you understand each symbol, you understand the proof. Not necessarily. You also need to understand the omitted transitions between lines.
\end{trapbox}

\begin{labbox}
Take five proofs from earlier chapters. For each, write an expanded version and a compressed version. Then write one sentence naming what was compressed.
\end{labbox}

\subsection{Logic review}

Core definitions:

\begin{itemize}
\item A statement is a sentence that is true or false.
\item A predicate becomes a statement when variables receive values from specified domains.
\item A conditional $P\Rightarrow Q$ is false exactly when $P$ is true and $Q$ is false.
\item The contrapositive of $P\Rightarrow Q$ is $\neg Q\Rightarrow\neg P$.
\item The converse of $P\Rightarrow Q$ is $Q\Rightarrow P$.
\end{itemize}

Mastery questions:

\begin{enumerate}
\item Can you tell whether a sentence is a statement before trying to prove it?
\item Can you identify the domain of every variable?
\item Can you negate a conditional without saying ``if not $P$, then not $Q$''?
\item Can you explain why a statement and its contrapositive are equivalent?
\item Can you produce a counterexample to the converse of a true implication?
\end{enumerate}

\begin{labbox}
Write five true implications from elementary number theory. For each one, write the converse and decide whether it is true. If false, give a counterexample.
\end{labbox}

\subsection{Quantifier review}

Quantifiers are not decoration. They are instructions for proof and disproof.

\begin{longtable}{p{0.25\linewidth}p{0.3\linewidth}p{0.35\linewidth}}
\toprule
Statement form & To prove it & To disprove it \\
\midrule
$\forall x\in A,\ P(x)$ & Let arbitrary $x\in A$ and prove $P(x)$ & Find one $x\in A$ with $\neg P(x)$ \\
$\exists x\in A,\ P(x)$ & Exhibit a particular $x\in A$ with $P(x)$ & Prove every $x\in A$ fails $P(x)$ \\
$\forall x\exists y,\ R(x,y)$ & Let arbitrary $x$, then construct $y$ depending on $x$ & Find one $x$ for which no $y$ works \\
$\exists y\forall x,\ R(x,y)$ & Produce one $y$ that works for all $x$ & Show every candidate $y$ fails for some $x$ \\
\bottomrule
\end{longtable}

Drill:

\begin{enumerate}
\item Prove or disprove: $\forall x\in\R,\exists y\in\R,\ x+y=0$.
\item Prove or disprove: $\exists y\in\R,\forall x\in\R,\ x+y=0$.
\item Prove or disprove: $\forall n\in\Z,\exists m\in\Z,\ n< m$.
\item Prove or disprove: $\exists m\in\Z,\forall n\in\Z,\ n<m$.
\item Translate the previous two into everyday English.
\end{enumerate}

\subsection{Definition review}

A definition is a contract. It tells you what you may use and what you must prove.

\begin{formalbox}
To prove an object has property $P$, expand the definition of $P$ and satisfy each clause. To use the fact that an object has property $P$, expand the definition and extract its clauses.
\end{formalbox}

Example: divisibility.

\[
a\mid b\quad\text{means}\quad \exists k\in\Z\text{ such that }b=ak.
\]

To prove $a\mid b$, produce $k$. To use $a\mid b$, introduce $k$ with $b=ak$.

Example: injective.

\[
f\text{ injective}\quad\text{means}\quad f(x_1)=f(x_2)\Rightarrow x_1=x_2.
\]

To prove injective, assume equal outputs. To use injective, first obtain equal outputs.

Example: open set in a metric space.

\[
U\text{ open}\quad\text{means}\quad \forall x\in U,\exists r>0,\ B_r(x)\subseteq U.
\]

To prove open, start with arbitrary point and construct a radius. To use open, feed it a point and receive a radius.

\subsection{Direct proof review}

Direct proof is the default when the hypothesis naturally opens into the ingredients needed by the conclusion.

Claim: If $a,b,c\in\Z$, $a\mid b$, and $a\mid c$, then $a\mid(2b-3c)$.

\begin{proof}
Since $a\mid b$, there exists $m\in\Z$ such that $b=am$. Since $a\mid c$, there exists $n\in\Z$ such that $c=an$. Then
\[
2b-3c=2(am)-3(an)=a(2m-3n).
\]
Since $2m-3n\in\Z$, we have $a\mid(2b-3c)$.
\end{proof}

Commentary: The proof is almost mechanical because the conclusion asks for a multiple of $a$, and the hypotheses already express $b$ and $c$ as multiples of $a$.

\subsection{Contrapositive review}

Contrapositive proof is useful when the conclusion is hard to show directly but its negation has a clean structure.

Claim: If $n^2$ is divisible by $3$, then $n$ is divisible by $3$.

\begin{proof}
We prove the contrapositive. Suppose $n$ is not divisible by $3$. Then $n$ has remainder $1$ or $2$ when divided by $3$. If $n=3q+1$, then
\[
n^2=(3q+1)^2=9q^2+6q+1=3(3q^2+2q)+1,
\]
so $n^2$ is not divisible by $3$. If $n=3q+2$, then
\[
n^2=(3q+2)^2=9q^2+12q+4=3(3q^2+4q+1)+1,
\]
so again $n^2$ is not divisible by $3$. Therefore, if $n^2$ is divisible by $3$, then $n$ is divisible by $3$.
\end{proof}

\begin{intuitionbox}
The contrapositive turns ``a hidden factor in $n^2$ forces a hidden factor in $n$'' into a visible remainder calculation.
\end{intuitionbox}

\subsection{Contradiction review}

Proof by contradiction is powerful, but it should end in a named impossibility.

Weak ending:

\begin{quote}
This is a contradiction.
\end{quote}

Better ending:

\begin{quote}
This contradicts the fact that $p$ and $q$ have no common divisor greater than $1$.
\end{quote}

Classic example:

\begin{proof}
Suppose, for contradiction, that $\sqrt2$ is rational. Then $\sqrt2=p/q$ for integers $p,q$ with $q\neq0$ and with the fraction in lowest terms. Squaring gives
\[
2=\frac{p^2}{q^2},
\]
so
\[
p^2=2q^2.
\]
Thus $p^2$ is even, so $p$ is even. Write $p=2r$. Then
\[
4r^2=2q^2,
\]
so
\[
q^2=2r^2.
\]
Thus $q^2$ is even, so $q$ is even. Hence both $p$ and $q$ are even, contradicting that $p/q$ was in lowest terms. Therefore $\sqrt2$ is irrational.
\end{proof}

Commentary: The proof uses a lemma twice: if a square is even, then the original integer is even. In a polished proof, cite it. In a beginner portfolio, prove it before using it.

\subsection{Cases review}

Cases are useful only when they are exhaustive and each case reaches the same target.

Claim: For every integer $n$, $n^2+n$ is even.

\begin{proof}
Let $n\in\Z$. There are two cases.

If $n$ is even, then $n=2k$ for some $k\in\Z$. Then
\[
n^2+n=n(n+1)=2k(n+1),
\]
which is even.

If $n$ is odd, then $n=2k+1$ for some $k\in\Z$. Then
\[
n+1=2k+2=2(k+1),
\]
so
\[
n^2+n=n(n+1)=2(n(k+1)),
\]
which is even.

In both cases, $n^2+n$ is even.
\end{proof}

Debugging question: where did we use that every integer is either even or odd? It appears in the phrase ``There are two cases.'' A fully formal proof would state the parity dichotomy as a known lemma.

\subsection{Set proof review}

Set proofs are logic proofs wearing membership notation.

Claim:
\[
A\setminus(B\cup C)=(A\setminus B)\cap(A\setminus C).
\]

\begin{proof}
Let $x$ be arbitrary. Then $x\in A\setminus(B\cup C)$ if and only if $x\in A$ and $x\notin B\cup C$. This is equivalent to $x\in A$ and not $(x\in B\text{ or }x\in C)$. By De Morgan's law, this is equivalent to
\[
x\in A,\qquad x\notin B,\qquad x\notin C.
\]
That is equivalent to $x\in A\setminus B$ and $x\in A\setminus C$, which is equivalent to
\[
x\in(A\setminus B)\cap(A\setminus C).
\]
Therefore the sets are equal.
\end{proof}

\subsection{Function proof review}

For functions, most beginner errors are starting with the wrong object.

\begin{longtable}{p{0.25\linewidth}p{0.33\linewidth}p{0.32\linewidth}}
\toprule
Target & Correct first move & Common wrong first move \\
\midrule
Prove $f$ injective & Let $x_1,x_2$ be in the domain and assume $f(x_1)=f(x_2)$ & Let $y$ be in the codomain \\
Prove $f$ surjective & Let $y$ be arbitrary in the codomain & Let $x$ be arbitrary in the domain \\
Prove $f^{-1}(S)\subseteq f^{-1}(T)$ & Let $x\in f^{-1}(S)$ & Assume $f$ has an inverse \\
Prove $f(A)\subseteq B$ & Let $y\in f(A)$ or let $x\in A$ and show $f(x)\in B$ & Treat $A$ as a number \\
\bottomrule
\end{longtable}

Drill:

\begin{enumerate}
\item Let $f:\R\to\R$ be defined by $f(x)=3x-5$. Prove $f$ is injective.
\item Let $f:\R\to\R$ be defined by $f(x)=3x-5$. Prove $f$ is surjective.
\item Let $f:\R\to\R$ be defined by $f(x)=x^2$. Prove $f$ is not injective.
\item Let $f:\R\to\R$ be defined by $f(x)=x^2$. Prove $f$ is not surjective.
\item Let $f:A\to B$. Prove that if $S\subseteq T\subseteq B$, then $f^{-1}(S)\subseteq f^{-1}(T)$.
\end{enumerate}

\subsection{Relation review}

Relations require property-by-property discipline. Do not compress reflexive, symmetric, and transitive into one sentence unless the relation is already routine for your reader.

Claim: On $\R$, define $x\sim y$ if $x-y\in\Z$. Then $\sim$ is an equivalence relation.

\begin{proof}
Reflexive: For any $x\in\R$, $x-x=0\in\Z$, so $x\sim x$.

Symmetric: Suppose $x\sim y$. Then $x-y\in\Z$. Since integers are closed under additive inverses,
\[
y-x=-(x-y)\in\Z.
\]
Thus $y\sim x$.

Transitive: Suppose $x\sim y$ and $y\sim z$. Then $x-y\in\Z$ and $y-z\in\Z$. Since integers are closed under addition,
\[
x-z=(x-y)+(y-z)\in\Z.
\]
Thus $x\sim z$.

Therefore $\sim$ is an equivalence relation.
\end{proof}

Later view: this relation identifies real numbers that differ by an integer. The quotient $\R/\Z$ can be pictured as a circle. This is a first glimpse of how equivalence relations create new spaces.

\subsection{Induction review}

An induction proof must make the induction statement clear. The statement is not ``the formula.'' It is a predicate $P(n)$.

Claim:
\[
1+3+5+\cdots+(2n-1)=n^2
\]
for all $n\in\N$.

\begin{proof}
Let $P(n)$ be the statement
\[
1+3+\cdots+(2n-1)=n^2.
\]
Base case: For $n=1$, the left side is $1$ and the right side is $1^2=1$, so $P(1)$ is true.

Induction step: Assume $P(k)$ is true for some $k\in\N$, so
\[
1+3+\cdots+(2k-1)=k^2.
\]
Then
\[
1+3+\cdots+(2k-1)+(2(k+1)-1)=k^2+(2k+1).
\]
Since
\[
k^2+2k+1=(k+1)^2,
\]
we get
\[
1+3+\cdots+(2(k+1)-1)=(k+1)^2.
\]
Thus $P(k+1)$ is true. By induction, $P(n)$ is true for all $n\in\N$.
\end{proof}

\subsection{Spaced repetition prompts}

Use these prompts on a rotating schedule. Answer from memory first, then check notes.

\begin{longtable}{p{0.15\linewidth}p{0.75\linewidth}}
\toprule
Day & Prompts \\
\midrule
1 & Define statement, predicate, implication, converse, contrapositive. Give one example of each. \\
2 & Negate three quantified statements: one $\forall$, one $\exists$, one $\forall\exists$. \\
3 & Prove a parity claim from definitions. Then write one bad proof of the same claim and repair it. \\
4 & Prove one subset relation and one set equality. Name each logical law used. \\
5 & Prove one injectivity claim and one surjectivity claim. Start with the correct arbitrary object. \\
6 & Verify one equivalence relation. Write the equivalence class of a sample element. \\
7 & Write one induction proof. Underline the base case, hypothesis, and induction step. \\
14 & Redo the hardest proof from week 1 without looking. Compare to the old version. \\
21 & Explain quantifier order using your own example. Include a false reversed version. \\
30 & Take a mini-exam: one logic question, one set proof, one function proof, one induction proof. \\
\bottomrule
\end{longtable}

\subsection{Mastery checklist}

Mark each item as solid, shaky, or missing.

\begin{enumerate}
\item I can write the definition of even, odd, divides, rational, injective, surjective, equivalence relation, partial order, and open set.
\item I can prove simple parity claims without relying on slogans.
\item I can negate nested quantified statements.
\item I can choose between direct proof, contrapositive, contradiction, and cases.
\item I can prove set equalities by arbitrary elements.
\item I can avoid confusing image and inverse image.
\item I can produce counterexamples using small finite sets or simple functions.
\item I can write clean induction proofs.
\item I can find the first bad line in a flawed proof.
\item I can explain a proof aloud in ordinary mathematical English.
\end{enumerate}

\begin{focusbox}
If a checklist item is shaky, assign it a proof task. Do not assign it more reading. Reading helps after production reveals the exact gap.
\end{focusbox}

\section{Portfolio Assignments}
\chapterline

A proof portfolio is a curated record of competence. It should not contain every exercise you have ever attempted. It should contain carefully revised proofs that show you can handle the core forms of early pure mathematics.

\subsection{Portfolio rules}

For each portfolio proof:

\begin{enumerate}
\item State the theorem before proving it.
\item List definitions you expect to use.
\item Write a proof skeleton.
\item Write the full proof.
\item Add a short commentary paragraph.
\item Add one possible mistake and explain how to avoid it.
\end{enumerate}

The commentary matters. It forces you to notice the architecture instead of treating the proof as a string of symbols.

\subsection{Portfolio A: arithmetic and definitions}

Write polished proofs of the following.

\begin{enumerate}
\item The sum of two even integers is even.
\item The sum of two odd integers is even.
\item The product of two odd integers is odd.
\item If $a\mid b$ and $a\mid c$, then $a\mid(b+c)$.
\item If $a\mid b$ and $b\mid c$, then $a\mid c$.
\item If $x,y\in\Q$, then $x-y\in\Q$.
\item If $x,y\in\Q$, then $xy\in\Q$.
\item If $x\in\Q$ and $x\neq0$, then $1/x\in\Q$.
\end{enumerate}

Required commentary:

\begin{quote}
For each proof, identify the witness required by the conclusion. For example, if proving $a\mid(b+c)$, identify the integer $k$ such that $b+c=ak$.
\end{quote}

\subsection{Portfolio B: proof methods}

Write polished proofs of the following.

\begin{enumerate}
\item If $n^2$ is even, then $n$ is even.
\item If $n^2$ is odd, then $n$ is odd.
\item If $n^2$ is divisible by $3$, then $n$ is divisible by $3$.
\item There is no largest integer.
\item There is no smallest positive real number.
\item $\sqrt2$ is irrational.
\item There are infinitely many primes.
\item For every integer $n$, $n^2+n$ is even.
\end{enumerate}

Required commentary:

\begin{quote}
Label the proof method. If you use contradiction, state the exact contradiction. If you use contrapositive, write both the original statement and the contrapositive.
\end{quote}

\subsection{Portfolio C: sets}

Write polished proofs of the following.

\begin{enumerate}
\item $A\cap B\subseteq A$.
\item $A\subseteq A\cup B$.
\item $A\setminus B=A\cap B^c$.
\item $A\cap(B\cup C)=(A\cap B)\cup(A\cap C)$.
\item $A\setminus(B\cup C)=(A\setminus B)\cap(A\setminus C)$.
\item If $A\subseteq B$, then $A\cap C\subseteq B\cap C$.
\item If $A\subseteq B$, then $\pow(A)\subseteq\pow(B)$.
\item $A\subseteq B$ if and only if $B^c\subseteq A^c$.
\end{enumerate}

Required commentary:

\begin{quote}
For each proof, translate the set statement into a logical statement about an arbitrary element $x$.
\end{quote}

\subsection{Portfolio D: functions}

Write polished proofs or counterexamples.

\begin{enumerate}
\item If $f:A\to B$ and $g:B\to C$ are injective, then $g\circ f$ is injective.
\item If $f:A\to B$ and $g:B\to C$ are surjective, then $g\circ f$ is surjective.
\item If $f:A\to B$ is injective and $S,T\subseteq A$, then $f(S\cap T)=f(S)\cap f(T)$.
\item Give a counterexample showing that the previous equality can fail when $f$ is not injective.
\item Prove $f^{-1}(S\cap T)=f^{-1}(S)\cap f^{-1}(T)$.
\item Prove $f^{-1}(S\cup T)=f^{-1}(S)\cup f^{-1}(T)$.
\item Prove that if $S\subseteq T$, then $f^{-1}(S)\subseteq f^{-1}(T)$.
\item For $f:\R\to\R$ given by $f(x)=x^2$, compute $f([-2,1])$ and $f^{-1}([0,4])$.
\end{enumerate}

Required commentary:

\begin{quote}
For each function proof, state whether you start with an element of the domain, codomain, image, or inverse image. Explain why.
\end{quote}

\subsection{Portfolio E: relations and induction}

Write polished proofs of the following.

\begin{enumerate}
\item Congruence modulo $m$ is an equivalence relation on $\Z$.
\item The relation on $\R$ defined by $x\sim y$ if $x-y\in\Z$ is an equivalence relation.
\item The equivalence classes of congruence modulo $4$ partition $\Z$.
\item Subset inclusion is a partial order on $\pow(X)$.
\item Divisibility is a partial order on $\N$.
\item Prove by induction that $1+2+\cdots+n=n(n+1)/2$.
\item Prove by induction that $1+3+\cdots+(2n-1)=n^2$.
\item Prove by induction that $2^n\geq n+1$ for all $n\geq1$.
\end{enumerate}

Required commentary:

\begin{quote}
For relation proofs, describe each property separately. For induction proofs, identify $P(n)$ before writing the proof.
\end{quote}

\section{Guided Solutions and Commentary}
\chapterline

Try each problem before reading its solution. The point of a guided solution is not just to show the answer. It is to show how the proof was found, where the definitions enter, and what a reader should notice.

\subsection{Solution 1: sum of two even integers}

Claim: If $a$ and $b$ are even integers, then $a+b$ is even.

Plan: Open the definition of even for $a$ and for $b$. The conclusion asks us to express $a+b$ as $2$ times an integer.

\begin{proof}
Let $a,b\in\Z$ be even. Then there exist $m,n\in\Z$ such that
\[
a=2m,\qquad b=2n.
\]
Therefore
\[
a+b=2m+2n=2(m+n).
\]
Since $m+n\in\Z$, the number $a+b$ is even.
\end{proof}

Commentary: The witness for evenness of $a+b$ is $m+n$.

\subsection{Solution 2: product of odd integers}

Claim: If $a$ and $b$ are odd integers, then $ab$ is odd.

Plan: Odd means $2k+1$. Multiply the forms and factor the result into $2(\text{integer})+1$.

\begin{proof}
Let $a,b\in\Z$ be odd. Then $a=2m+1$ and $b=2n+1$ for some $m,n\in\Z$. Hence
\[
ab=(2m+1)(2n+1)=4mn+2m+2n+1=2(2mn+m+n)+1.
\]
Since $2mn+m+n\in\Z$, $ab$ is odd.
\end{proof}

Commentary: Do not stop at $4mn+2m+2n+1$. The final form must visibly match the definition of odd.

\subsection{Solution 3: divisibility is transitive}

Claim: If $a\mid b$ and $b\mid c$, then $a\mid c$.

\begin{proof}
Suppose $a\mid b$ and $b\mid c$. Then $b=am$ and $c=bn$ for some $m,n\in\Z$. Therefore
\[
c=bn=(am)n=a(mn).
\]
Since $mn\in\Z$, $a\mid c$.
\end{proof}

Commentary: Divisibility proofs are witness proofs. The witness for $a\mid c$ is $mn$.

\subsection{Solution 4: rational subtraction}

Claim: If $x,y\in\Q$, then $x-y\in\Q$.

\begin{proof}
Let $x,y\in\Q$. Then there exist integers $a,b,c,d$ with $b\neq0$ and $d\neq0$ such that
\[
x=\frac ab,\qquad y=\frac cd.
\]
Thus
\[
x-y=\frac ab-\frac cd=\frac{ad-bc}{bd}.
\]
Since $ad-bc\in\Z$ and $bd\in\Z$ with $bd\neq0$, the number $x-y$ is rational.
\end{proof}

Commentary: The denominator condition must be preserved. Since $b$ and $d$ are nonzero integers, $bd$ is nonzero.

\subsection{Solution 5: no largest integer}

Claim: There is no largest integer.

\begin{proof}
Suppose, for contradiction, that there is a largest integer $N$. Then $N+1$ is an integer and $N+1>N$. This contradicts the assumption that $N$ is largest. Therefore there is no largest integer.
\end{proof}

Commentary: The contradiction is not emotional. It is the existence of an integer larger than the alleged largest integer.

\subsection{Solution 6: a set containment}

Claim: If $A\subseteq B$, then $A\cap C\subseteq B\cap C$.

\begin{proof}
Suppose $A\subseteq B$. Let $x\in A\cap C$. Then $x\in A$ and $x\in C$. Since $A\subseteq B$, $x\in B$. Therefore $x\in B$ and $x\in C$, so $x\in B\cap C$. Hence $A\cap C\subseteq B\cap C$.
\end{proof}

Commentary: The arbitrary element begins in the smaller set $A\cap C$, because the target is a subset claim.

\subsection{Solution 7: power sets}

Claim: If $A\subseteq B$, then $\pow(A)\subseteq\pow(B)$.

\begin{proof}
Suppose $A\subseteq B$. Let $S\in\pow(A)$. Then $S\subseteq A$. Since $A\subseteq B$, it follows that $S\subseteq B$. Thus $S\in\pow(B)$. Therefore $\pow(A)\subseteq\pow(B)$.
\end{proof}

Commentary: The elements of $\pow(A)$ are sets. The proof respects that type.

\subsection{Solution 8: inverse image of a union}

Claim:
\[
f^{-1}(S\cup T)=f^{-1}(S)\cup f^{-1}(T).
\]

\begin{proof}
Let $x$ be arbitrary. Then
\[
x\in f^{-1}(S\cup T)
\]
if and only if $f(x)\in S\cup T$. This holds if and only if $f(x)\in S$ or $f(x)\in T$. That is equivalent to $x\in f^{-1}(S)$ or $x\in f^{-1}(T)$, which is equivalent to
\[
x\in f^{-1}(S)\cup f^{-1}(T).
\]
Therefore the two sets are equal.
\end{proof}

Commentary: Inverse image proofs are usually clean because membership translates directly through $f$.

\subsection{Solution 9: image of intersection can fail}

Claim: It is not always true that
\[
f(S\cap T)=f(S)\cap f(T).
\]

Counterexample: Let $f:\R\to\R$ be $f(x)=x^2$, let $S=\{-1\}$, and let $T=\{1\}$. Then
\[
S\cap T=\emptyset,
\]
so
\[
f(S\cap T)=f(\emptyset)=\emptyset.
\]
But
\[
f(S)=\{1\},\qquad f(T)=\{1\},
\]
so
\[
f(S)\cap f(T)=\{1\}.
\]
Thus $f(S\cap T)\neq f(S)\cap f(T)$.

Commentary: The failure happens because two different inputs have the same output. This is exactly the kind of collision injectivity forbids.

\subsection{Solution 10: composition of injective functions}

Claim: If $f:A\to B$ and $g:B\to C$ are injective, then $g\circ f:A\to C$ is injective.

\begin{proof}
Let $a_1,a_2\in A$ and suppose $(g\circ f)(a_1)=(g\circ f)(a_2)$. Then
\[
g(f(a_1))=g(f(a_2)).
\]
Since $g$ is injective, $f(a_1)=f(a_2)$. Since $f$ is injective, $a_1=a_2$. Therefore $g\circ f$ is injective.
\end{proof}

Commentary: Injectivity is used backward through the composition: first $g$, then $f$.

\subsection{Solution 11: composition of surjective functions}

Claim: If $f:A\to B$ and $g:B\to C$ are surjective, then $g\circ f:A\to C$ is surjective.

\begin{proof}
Let $c\in C$ be arbitrary. Since $g$ is surjective, there exists $b\in B$ such that $g(b)=c$. Since $f$ is surjective, there exists $a\in A$ such that $f(a)=b$. Then
\[
(g\circ f)(a)=g(f(a))=g(b)=c.
\]
Thus $g\circ f$ is surjective.
\end{proof}

Commentary: Surjectivity is proved by chasing a target backward from the codomain.

\subsection{Solution 12: congruence modulo m}

Claim: For $m\geq2$, congruence modulo $m$ is an equivalence relation on $\Z$.

\begin{proof}
For $a,b\in\Z$, define $a\sim b$ if $m\mid(a-b)$.

Reflexive: Since $a-a=0=m\cdot0$, we have $m\mid(a-a)$, so $a\sim a$.

Symmetric: If $a\sim b$, then $a-b=mk$ for some $k\in\Z$. Hence $b-a=m(-k)$, so $b\sim a$.

Transitive: If $a\sim b$ and $b\sim c$, then $a-b=mk$ and $b-c=m\ell$ for some $k,\ell\in\Z$. Therefore
\[
a-c=(a-b)+(b-c)=m(k+\ell),
\]
so $a\sim c$.

Thus $\sim$ is an equivalence relation.
\end{proof}

Commentary: The same divisibility definition powers all three parts.

\subsection{Solution 13: subset inclusion is a partial order}

Claim: $\subseteq$ is a partial order on $\pow(X)$.

\begin{proof}
Reflexive: If $A\in\pow(X)$, then every element of $A$ is in $A$, so $A\subseteq A$.

Antisymmetric: Suppose $A,B\in\pow(X)$, $A\subseteq B$, and $B\subseteq A$. Then $A$ and $B$ have the same elements, so $A=B$.

Transitive: Suppose $A\subseteq B$ and $B\subseteq C$. If $x\in A$, then $x\in B$, and since $B\subseteq C$, $x\in C$. Hence $A\subseteq C$.

Therefore $\subseteq$ is a partial order on $\pow(X)$.
\end{proof}

Commentary: Antisymmetry is not symmetry. Symmetry would say $A\subseteq B$ implies $B\subseteq A$, which is false.

\subsection{Solution 14: induction for powers of two}

Claim: For all $n\geq1$,
\[
2^n\geq n+1.
\]

\begin{proof}
Base case: If $n=1$, then $2^1=2=1+1$, so the claim holds.

Induction step: Suppose $2^k\geq k+1$ for some $k\geq1$. Then
\[
2^{k+1}=2\cdot2^k\geq2(k+1)=2k+2.
\]
Since $2k+2\geq k+2$ for $k\geq0$, we have
\[
2^{k+1}\geq k+2.
\]
This is the desired statement for $k+1$. Therefore $2^n\geq n+1$ for all $n\geq1$.
\end{proof}

Commentary: The target for the induction step is $2^{k+1}\geq(k+1)+1=k+2$.

\subsection{Solution 15: finite intersection of open sets}

Claim: If $U$ and $V$ are open subsets of $\R$, then $U\cap V$ is open.

\begin{proof}
Let $x\in U\cap V$. Then $x\in U$ and $x\in V$. Since $U$ is open, there exists $r_1>0$ such that $(x-r_1,x+r_1)\subseteq U$. Since $V$ is open, there exists $r_2>0$ such that $(x-r_2,x+r_2)\subseteq V$.

Let $r=\min\{r_1,r_2\}$. Then $r>0$. If $y\in(x-r,x+r)$, then $|y-x|<r\leq r_1$, so $y\in U$, and $|y-x|<r\leq r_2$, so $y\in V$. Hence $y\in U\cap V$.

Thus every $x\in U\cap V$ has a positive-radius interval contained in $U\cap V$, so $U\cap V$ is open.
\end{proof}

Commentary: The smaller radius is the witness that works for both open sets.

\section{Proof Repair Workshops}
\chapterline

The previous debugging section gave individual repairs. This section gives workshops: clusters of problems that share the same kind of failure. The goal is pattern recognition. When you see a certain kind of broken proof, you should begin to know what kind of repair is likely.

\begin{focusbox}
Do not memorize repairs as scripts. Learn the pressure points: quantifiers, witnesses, definitions, cases, and direction of implication.
\end{focusbox}

\section{Reading Real Mathematics}
\chapterline

Independent mathematical literacy begins when you can unpack compressed prose.
Real books and papers often omit routine steps, change notation quickly, and
expect you to infer which theorem is being used. This section trains that skill
directly.

\begin{marginnotebox}
\motifhiddenquantifier\ Hidden quantifier. Real mathematical prose often hides
``for every'' and ``there exists'' inside phrases like ``choose,'' ``fix,''
``locally,'' ``sufficiently large,'' and ``for all small enough.''
\end{marginnotebox}

\subsection{Compressed excerpt 1: set containment}

Compressed text:

\begin{quote}
Since $A\subseteq B$, it follows that $\pow(A)\subseteq\pow(B)$.
\end{quote}

Expanded reading:

\begin{enumerate}
\item To prove $\pow(A)\subseteq\pow(B)$, let $S\in\pow(A)$.
\item Then $S\subseteq A$.
\item Since $A\subseteq B$, every element of $S$ is in $B$.
\item Hence $S\subseteq B$.
\item Therefore $S\in\pow(B)$.
\end{enumerate}

Hidden assumption: $\pow(A)$ means the set of all subsets of $A$.

\subsection{Compressed excerpt 2: continuity by preimages}

Compressed text:

\begin{quote}
The composition is continuous because preimages commute with composition.
\end{quote}

Expanded reading:

\begin{enumerate}
\item Let $f:X\to Y$ and $g:Y\to Z$ be continuous.
\item To prove $g\circ f$ continuous, take an open set $W\subseteq Z$.
\item Since $g$ is continuous, $g^{-1}(W)$ is open in $Y$.
\item Since $f$ is continuous, $f^{-1}(g^{-1}(W))$ is open in $X$.
\item But $f^{-1}(g^{-1}(W))=(g\circ f)^{-1}(W)$.
\item Therefore $(g\circ f)^{-1}(W)$ is open in $X$.
\end{enumerate}

Hidden theorem: inverse images under composition satisfy
\[
(g\circ f)^{-1}(W)=f^{-1}(g^{-1}(W)).
\]

\subsection{Compressed excerpt 3: compactness}

Compressed text:

\begin{quote}
By compactness, pass to a finite subcover.
\end{quote}

Expanded reading:

\begin{enumerate}
\item There must be an open cover somewhere in the argument.
\item Compactness applies only after you verify the covered set is compact.
\item A finite subcover means finitely many sets from the original cover still cover the same compact set.
\item The proof should name those finitely many indices.
\end{enumerate}

Hidden assumption: the cover consists of open sets in the correct topology.

\subsection{Hidden-assumption detection}

When reading a proof, ask:

\begin{enumerate}
\item Was a variable introduced with a domain?
\item Was a theorem applied with all hypotheses checked?
\item Was a map assumed continuous, injective, surjective, or well-defined?
\item Was a set open or closed in the correct ambient space?
\item Was division by a possibly zero quantity avoided?
\item Was a finite subcover taken only after compactness was established?
\end{enumerate}

\begin{labbox}
Notation expansion drill. Expand each phrase into three to six explicit lines:
\begin{enumerate}
\item ``By injectivity, the result follows.''
\item ``The reverse inclusion is similar.''
\item ``Choose $N$ sufficiently large.''
\item ``This is well-defined.''
\item ``The quotient map is continuous by definition.''
\end{enumerate}
\end{labbox}

\subsection{Theorem Dependency Tracing}

Advanced reading often fails because a proof silently leans on earlier theorems.
You may understand every line locally and still miss the dependency chain.

\begin{formalbox}
Dependency trace template:
\begin{enumerate}
\item Name the theorem being proved.
\item List each theorem, lemma, or definition invoked.
\item For each invocation, check its hypotheses.
\item Mark whether the dependency is local, from the current chapter, or from a previous field.
\item Ask what breaks if the dependency is unavailable.
\end{enumerate}
\end{formalbox}

Example:

\begin{quote}
Every continuous function $f:[a,b]\to\R$ attains a maximum.
\end{quote}

Dependency chain:

\begin{center}
\begin{tikzpicture}[
  node distance=0.9cm,
  dep/.style={draw, rounded corners, align=center, fill=gray!5, minimum width=4.3cm, minimum height=0.72cm},
  arrow/.style={-{Latex}, thick}
]
\node[dep] (hb) {Heine-Borel:\\$[a,b]$ is compact};
\node[dep, below=of hb] (ci) {Continuous image\\of compact is compact};
\node[dep, below=of ci] (cr) {Compact subsets of $\R$\\are closed and bounded};
\node[dep, below=of cr] (sup) {Supremum property\\of $\R$};
\node[dep, below=of sup] (max) {$f$ attains a maximum};
\draw[arrow] (hb) -- (ci);
\draw[arrow] (ci) -- (cr);
\draw[arrow] (cr) -- (sup);
\draw[arrow] (sup) -- (max);
\end{tikzpicture}
\end{center}

This is why compactness matters. It is not merely a definition about covers; it
is a machine that feeds existence theorems.

\subsection{Surviving Abstraction Jumps}

Mathematical texts often move from examples to abstractions quickly:

\[
\text{integers modulo }n
\quad\leadsto\quad
\text{quotient groups}
\quad\leadsto\quad
\text{quotient spaces}
\quad\leadsto\quad
\text{moduli spaces}.
\]

When this happens, do not try to swallow the abstraction whole. Ask:

\begin{enumerate}
\item What was the original example?
\item What features of the example are being preserved?
\item What details are being deliberately forgotten?
\item What maps are allowed in the new setting?
\item What counts as sameness now?
\end{enumerate}

\begin{aestheticbox}
Abstraction feels powerful when it compresses many examples without erasing the
reason they mattered. Bad abstraction feels like fog. Good abstraction feels
like suddenly seeing the same architecture in several cities.
\end{aestheticbox}

\subsection{Paper-Style Proof Drill}

Here is a deliberately compressed proof in a paper-like style.

\begin{quote}
Let $q:X\to X/{\sim}$ be the quotient map. Since $U\subseteq X/{\sim}$ is open
iff $q^{-1}(U)$ is open, $q$ is continuous. Moreover, any map $f:X/{\sim}\to Y$
is continuous exactly when $f\circ q$ is continuous, provided $f$ is compatible
with the quotient structure.
\end{quote}

Unpack it:

\begin{enumerate}
\item What is the definition of the quotient topology?
\item Why does that definition immediately imply $q$ is continuous?
\item What does ``compatible with the quotient structure'' mean?
\item What hidden universal property is being hinted at?
\item Which direction of the continuity claim is immediate, and which requires the quotient topology?
\end{enumerate}

\begin{marginnotebox}
The phrase ``universal property'' is common in modern mathematics. It usually
means an object is characterized by how maps into or out of it behave. This is
category-theoretic thinking appearing before formal category theory.
\end{marginnotebox}

\subsection{Notation Fatigue Protocol}

When notation fatigue hits, do not keep rereading the same line at full speed.
Use a reset protocol.

\begin{enumerate}
\item Replace every symbol by its type: number, set, function, space, relation.
\item Rename the main objects in ordinary words.
\item Identify the active map, if there is one.
\item Write the target in plain English.
\item Reinsert the notation only after the sentence has meaning again.
\end{enumerate}

Example:
\[
q^{-1}(U)\text{ is open in }X.
\]

Types:

\begin{itemize}
\item $q:X\to X/{\sim}$ is a quotient map,
\item $U$ is a subset of the quotient space,
\item $q^{-1}(U)$ is the set of original points whose equivalence classes land in $U$.
\end{itemize}

Plain English:

\begin{quote}
The set of original points that map into $U$ is open upstairs.
\end{quote}

\begin{recoverybox}
If you understand the words but not the meaning, translate the sentence into
types. Meaning often returns when every symbol has a job again.
\end{recoverybox}

\subsection{A Gentle Recovery Chapter: One Proof, Three Readings}

Here is a small proof read three ways.

Claim:

\[
A\cap B\subseteq A.
\]

Beginner reading:

\begin{quote}
If $x$ is in $A\cap B$, then $x$ is in both $A$ and $B$. So $x$ is in $A$.
\end{quote}

Formal reading:

\begin{proof}
Let $x\in A\cap B$. By definition of intersection, $x\in A$ and $x\in B$.
Therefore $x\in A$. Hence $A\cap B\subseteq A$.
\end{proof}

Compressed reading:

\begin{quote}
Immediate from the definition of intersection.
\end{quote}

All three are the same mathematics at different compression levels.

\begin{aestheticbox}
Mathematical maturity is not always knowing more facts. Sometimes it is knowing
how much compression a situation can safely tolerate.
\end{aestheticbox}

\subsection{Workshop 1: quantifier inversion}

Quantifier inversion happens when a proof silently changes the order of choices. This is one of the most common errors in analysis, topology, and algebra.

Claim:
\[
\forall x\in\R,\exists y\in\R,\ y>x.
\]

Correct proof:

\begin{proof}
Let $x\in\R$ be arbitrary. Choose $y=x+1$. Then $y\in\R$ and $y>x$. Therefore for every real $x$, there exists a real $y$ with $y>x$.
\end{proof}

Bad proof of a different claim:

\begin{quote}
Choose $y=1000$. Then for every $x$, $y>x$.
\end{quote}

This attempts to prove
\[
\exists y\in\R,\forall x\in\R,\ y>x,
\]
which is false. No single real number is larger than every real number.

Repair habit: write the quantifier game.

\begin{itemize}
\item For $\forall x\exists y$, the opponent gives $x$ first; you may choose $y$ after seeing $x$.
\item For $\exists y\forall x$, you must choose $y$ first; it must survive every $x$.
\end{itemize}

Drill:

\begin{enumerate}
\item Prove $\forall x\in\R,\exists y\in\R,\ x+y=0$.
\item Disprove $\exists y\in\R,\forall x\in\R,\ x+y=0$.
\item Prove $\forall n\in\Z,\exists m\in\Z,\ m>n$.
\item Disprove $\exists m\in\Z,\forall n\in\Z,\ m>n$.
\end{enumerate}

\begin{walkbox}
When you hear ``for every,'' imagine the problem handing you an input. When you hear ``there exists,'' imagine you must produce something concrete.
\end{walkbox}

\subsection{Workshop 2: missing witnesses}

Existence claims require witnesses. Sometimes the witness is a number. Sometimes it is a function, a set, a radius, an index, or an equivalence class.

Claim: If $a\mid b$ and $a\mid c$, then $a\mid(5b+7c)$.

Bad proof:

\begin{quote}
Since $a$ divides both $b$ and $c$, it divides any combination of them.
\end{quote}

The idea is good, but the witness is missing.

\begin{proof}
Suppose $a\mid b$ and $a\mid c$. Then $b=am$ and $c=an$ for some $m,n\in\Z$. Hence
\[
5b+7c=5am+7an=a(5m+7n).
\]
Since $5m+7n\in\Z$, we have $a\mid(5b+7c)$.
\end{proof}

Witness inventory:

\begin{longtable}{p{0.28\linewidth}p{0.32\linewidth}p{0.3\linewidth}}
\toprule
Target & Needed witness & Typical source \\
\midrule
$a\mid b$ & $k\in\Z$ with $b=ak$ & Algebraic factoring \\
$x\in\Q$ & integers $p,q$ with $q\neq0$ and $x=p/q$ & Common denominator \\
$f$ surjective & preimage $a$ for arbitrary codomain element $b$ & Solve $f(a)=b$ \\
$U$ open & radius $r>0$ for arbitrary point $x\in U$ & Distance to boundary \\
$a_n\to L$ & index $N$ for arbitrary $\eps>0$ & Inequality scratch work \\
$\exists x\,P(x)$ & object $x$ satisfying $P$ & Construction or known theorem \\
\bottomrule
\end{longtable}

Drill:

\begin{enumerate}
\item In the proof that $x+y\in\Q$ for rational $x,y$, what is the witness numerator and denominator?
\item In the proof that $f(x)=2x+1$ is surjective from $\R$ to $\R$, what preimage do you choose for a given $y$?
\item In the proof that $(0,3)$ is open in $\R$, what radius works for a point $x\in(0,3)$?
\item In the proof that $7/n\to0$, what $N$ works for a given $\eps>0$?
\end{enumerate}

\subsection{Workshop 3: uniqueness}

Uniqueness proofs have a standard shape:

\begin{enumerate}
\item Suppose there are two objects with the desired property.
\item Use the defining property of both.
\item Prove the two objects are equal.
\end{enumerate}

Claim: If a function has an inverse, then the inverse is unique.

Bad proof:

\begin{quote}
The inverse is unique because a function can only be undone one way.
\end{quote}

This is intuition, not proof.

\begin{proof}
Let $f:A\to B$. Suppose $g:B\to A$ and $h:B\to A$ are both inverses of $f$. Then
\[
g\circ f=\operatorname{id}_A,\quad f\circ g=\operatorname{id}_B,
\]
and
\[
h\circ f=\operatorname{id}_A,\quad f\circ h=\operatorname{id}_B.
\]
We show $g=h$. Let $b\in B$. Since $h$ is an inverse of $f$, we have $f(h(b))=b$. Therefore
\[
g(b)=g(f(h(b)))=(g\circ f)(h(b))=\operatorname{id}_A(h(b))=h(b).
\]
Since this holds for every $b\in B$, $g=h$.
\end{proof}

Commentary: The proof does not rely on a picture of undoing. It uses the equations that define an inverse.

Drill:

\begin{enumerate}
\item Prove that an identity element in a group is unique.
\item Prove that inverses in a group are unique.
\item Prove that if a real number has an additive inverse, it is unique.
\item Prove that if a maximum element of a subset of $\R$ exists, it is unique.
\end{enumerate}

\subsection{Workshop 4: counterexample construction}

To disprove a universal claim, one counterexample is enough. But a useful counterexample is not random; it targets the claim's weak point.

False claim:
\[
f(S\cap T)=f(S)\cap f(T)
\]
for every function $f$ and all subsets $S,T$ of the domain.

How to search:

\begin{itemize}
\item The left side uses points that are in both $S$ and $T$.
\item The right side can contain an output reached from one point in $S$ and a different point in $T$.
\item So make $S$ and $T$ disjoint, but make their images overlap.
\end{itemize}

Counterexample:

\[
f:\R\to\R,\quad f(x)=x^2,\quad S=\{-1\},\quad T=\{1\}.
\]

Then $S\cap T=\emptyset$, so $f(S\cap T)=\emptyset$, but $f(S)\cap f(T)=\{1\}$.

False claim:

\[
\text{If }f\circ g\text{ is injective, then }f\text{ is injective.}
\]

Search:

\begin{itemize}
\item To make $f$ not injective, it needs two different inputs with the same output.
\item But $g$ might never hit both of those inputs.
\item Choose $g$ with small image.
\end{itemize}

Counterexample: Let $g:\{0\}\to\{0,1\}$ be $g(0)=0$, and let $f:\{0,1\}\to\{0\}$ be the constant function. Then $f$ is not injective, but $f\circ g:\{0\}\to\{0\}$ is injective because its domain has one element.

\begin{trapbox}
Counterexamples often live in tiny sets. Do not search first among complicated objects when a two-element set will expose the failure.
\end{trapbox}

Drill:

\begin{enumerate}
\item Disprove: If $f\circ g$ is surjective, then $g$ is surjective.
\item Disprove: If $A\cup C=B\cup C$, then $A=B$.
\item Disprove: If $A\cap C=B\cap C$, then $A=B$.
\item Disprove: Every relation that is symmetric and transitive is reflexive.
\item Disprove: Every continuous function is differentiable.
\end{enumerate}

\subsection{Workshop 5: hidden division by zero}

Algebraic proofs often fail because they divide by a quantity that might be zero.

False proof that all numbers are equal:

\begin{quote}
Let $a=b$. Then $a^2=ab$. Subtract $b^2$:
\[
a^2-b^2=ab-b^2.
\]
Factor:
\[
(a-b)(a+b)=b(a-b).
\]
Cancel $a-b$:
\[
a+b=b.
\]
Since $a=b$, this gives $2b=b$, so $b=0$. Therefore all equal numbers are zero.
\end{quote}

The first bad line is the cancellation of $a-b$. Since $a=b$, $a-b=0$. Dividing or canceling by zero is invalid.

Repair habit:

\begin{itemize}
\item Before dividing by an expression, prove it is nonzero.
\item If it may be zero, split into cases.
\item If cancellation occurs in an abstract structure, check that the cancellation law is available.
\end{itemize}

Drill:

\begin{enumerate}
\item Identify the invalid step in the classic false proof that $1=2$.
\item Prove: If $ab=ac$ and $a\neq0$ in $\R$, then $b=c$.
\item Explain why the condition $a\neq0$ is necessary.
\item In a group, prove cancellation without using division.
\end{enumerate}

\subsection{Workshop 6: epsilon repairs}

Epsilon proofs are intimidating because the proof is written forward but discovered backward.

Claim:
\[
\frac{2n+3}{n}\to2.
\]

Scratch work:

\[
\left|\frac{2n+3}{n}-2\right|=\frac3n.
\]

We want $\frac3n<\eps$, so it is enough to choose $n>3/\eps$.

Final proof:

\begin{proof}
Let $\eps>0$. Choose $N\in\N$ such that $N>3/\eps$. If $n\geq N$, then
\[
\left|\frac{2n+3}{n}-2\right|=\frac3n\leq\frac3N<\eps.
\]
Therefore $(2n+3)/n\to2$.
\end{proof}

Bad proof:

\begin{quote}
Let $N>3/\eps$. Since $n$ is large, $3/n<\eps$.
\end{quote}

Diagnosis:

\begin{itemize}
\item It does not state $n\geq N$.
\item It says ``large'' instead of using the chosen threshold.
\item It does not show the absolute-value expression.
\end{itemize}

Drill:

\begin{enumerate}
\item Prove $5/n\to0$.
\item Prove $(n+1)/n\to1$.
\item Prove $(4n-7)/n\to4$.
\item Prove $1/(n+2)\to0$.
\end{enumerate}

\section{Final Exam: From Beginner to Velleman-Ready}
\chapterline

This exam is a diagnostic, not a performance ritual. Take it under realistic conditions: no notes for the main attempt, then a second pass with notes in a different color. The second pass is where learning becomes visible.

\begin{focusbox}
Scoring rule: a correct answer with unclear reasoning is not fully correct. A proof must show why the conclusion follows.
\end{focusbox}

\subsection{Instructions}

\begin{enumerate}
\item Allow three hours for the first attempt.
\item Write full sentences in proof problems.
\item State definitions when you use them.
\item For any problem you cannot solve, write the proof skeleton or explain the obstruction.
\item After finishing, debug your own work before checking solutions or notes.
\end{enumerate}

\subsection{Part 1: logic and language}

\begin{enumerate}
\item Define statement and predicate. Give two examples of each.
\item Explain why $x^2=9$ is not a statement until a domain and quantifier or value are supplied.
\item Translate into symbols: ``Every integer has a larger integer.''
\item Translate into English: $\exists M\in\R,\forall x\in S,\ x\leq M$.
\item Negate: $\forall x\in A,\exists y\in B,\ R(x,y)$.
\item Negate: $\exists N\in\N,\forall n\geq N,\ |a_n-L|<\eps$.
\item Give an example where $\forall x\exists y\,R(x,y)$ is true but $\exists y\forall x\,R(x,y)$ is false.
\item State the contrapositive and converse of: If $n^2$ is even, then $n$ is even.
\item Explain why a conditional with false hypothesis is true in classical logic.
\item Give a counterexample to the statement: If $n$ is odd, then $n$ is prime.
\end{enumerate}

\subsection{Part 2: definitions and arithmetic proofs}

\begin{enumerate}
\item Prove that the sum of two even integers is even.
\item Prove that the sum of two odd integers is even.
\item Prove that the product of two odd integers is odd.
\item Prove that if $a\mid b$ and $a\mid c$, then $a\mid(b-c)$.
\item Prove that if $a\mid b$ and $b\mid c$, then $a\mid c$.
\item Prove that if $x,y\in\Q$, then $x+y\in\Q$.
\item Prove that if $x,y\in\Q$, then $xy\in\Q$.
\item Prove that if $n^2$ is even, then $n$ is even.
\item Prove that if $n^2$ is divisible by $3$, then $n$ is divisible by $3$.
\item Prove that $\sqrt2$ is irrational.
\end{enumerate}

\subsection{Part 3: sets}

\begin{enumerate}
\item Prove $A\cap B\subseteq A$.
\item Prove $A\subseteq A\cup B$.
\item Prove $A\setminus B=A\cap B^c$.
\item Prove $A\cap(B\cup C)=(A\cap B)\cup(A\cap C)$.
\item Prove $A\setminus(B\cup C)=(A\setminus B)\cap(A\setminus C)$.
\item Prove that if $A\subseteq B$, then $A\cap C\subseteq B\cap C$.
\item Prove that if $A\subseteq B$, then $\pow(A)\subseteq\pow(B)$.
\item Prove $A\subseteq B$ if and only if $B^c\subseteq A^c$.
\item Give an example of two sets $A$ and $B$ such that $A\cup B=A$ but $B\neq A$.
\item Give an example showing that $A\setminus B$ need not equal $B\setminus A$.
\end{enumerate}

\subsection{Part 4: functions}

\begin{enumerate}
\item State the definitions of injective and surjective.
\item Prove that $f:\R\to\R$ given by $f(x)=3x-5$ is injective.
\item Prove that $f:\R\to\R$ given by $f(x)=3x-5$ is surjective.
\item Prove that if $f$ and $g$ are injective, then $g\circ f$ is injective.
\item Prove that if $f$ and $g$ are surjective, then $g\circ f$ is surjective.
\item Prove $f^{-1}(S\cap T)=f^{-1}(S)\cap f^{-1}(T)$.
\item Prove $f^{-1}(S\cup T)=f^{-1}(S)\cup f^{-1}(T)$.
\item Give a counterexample showing $f(S\cap T)=f(S)\cap f(T)$ can fail.
\item If $f:A\to B$ is injective and $S,T\subseteq A$, prove $f(S\cap T)=f(S)\cap f(T)$.
\item For $f:\R\to\R$ given by $f(x)=x^2$, compute $f^{-1}([1,4])$.
\end{enumerate}

\subsection{Part 5: relations}

\begin{enumerate}
\item Define reflexive, symmetric, transitive, antisymmetric, and total.
\item Prove congruence modulo $m$ is reflexive.
\item Prove congruence modulo $m$ is symmetric.
\item Prove congruence modulo $m$ is transitive.
\item Describe the equivalence classes modulo $5$.
\item Prove that if $\sim$ is an equivalence relation, then equivalence classes are either equal or disjoint.
\item Prove subset inclusion is a partial order on $\pow(X)$.
\item Prove divisibility is a partial order on $\N$.
\item Give an example of a relation that is reflexive and symmetric but not transitive.
\item Give an example of a relation that is transitive but not symmetric.
\end{enumerate}

\subsection{Part 6: induction}

\begin{enumerate}
\item State the principle of mathematical induction.
\item Prove by induction that $1+2+\cdots+n=n(n+1)/2$.
\item Prove by induction that $1+3+\cdots+(2n-1)=n^2$.
\item Prove by induction that $2^n\geq n+1$ for all $n\geq1$.
\item Prove by induction that $3\mid(4^n-1)$ for all $n\geq1$.
\item Prove by induction that $n! \geq 2^{n-1}$ for all $n\geq1$.
\item Explain the difference between ordinary induction and strong induction.
\item Use strong induction to prove that every integer $n\geq2$ has a prime divisor.
\item Identify the flaw: ``Assume $P(n)$ is true. Therefore $P(n+1)$ is true. By induction, $P(n)$ is true for all $n$.''
\item Write a proof skeleton for an induction proof before filling in algebra.
\end{enumerate}

\subsection{Part 7: proof debugging}

For each bad proof, identify the first bad line and repair it.

\begin{enumerate}
\item Claim: If $x,y\in\Q$, then $x+y\in\Q$. Bad proof: ``Write $x=a/b$ and $y=c/d$. Then $x+y=(a+c)/(b+d)$.''
\item Claim: If $n^2$ is even, then $n$ is even. Bad proof: ``Since $n^2=2k$, $n=\sqrt{2k}$, so $n$ is even.''
\item Claim: If $A\subseteq B$, then $\pow(A)\subseteq\pow(B)$. Bad proof: ``Power sets preserve size, so the result follows.''
\item Claim: If $f$ is surjective, then every element of the codomain is hit. Bad proof: ``Let $a$ be in the domain. Then $f(a)$ is in the codomain.''
\item Claim: $U\cap V$ is open if $U$ and $V$ are open. Bad proof: ``Open sets are closed under intersection, so it is open.''
\item Claim: Congruence modulo $m$ is symmetric. Bad proof: ``If $a$ and $b$ have the same remainder, then $b$ and $a$ have the same remainder.''
\item Claim: $f^{-1}(S\cap T)=f^{-1}(S)\cap f^{-1}(T)$. Bad proof: ``Apply $f^{-1}$ to both sides.''
\item Claim: There are infinitely many primes. Bad proof: ``The product of all primes plus $1$ is prime.''
\item Claim: If $A=B$, then $A\cup C=B\cup C$. Bad proof: ``This is obvious from equality.''
\item Claim: If $n$ is even, then $n^2$ is even. Bad proof: ``This is the converse of a true theorem.''
\end{enumerate}

\subsection{Part 8: readiness reflection}

Answer in paragraphs.

\begin{enumerate}
\item Which proof method feels most natural to you, and why?
\item Which proof method produces the most mistakes for you?
\item What is your most common proof error: quantifier, definition, witness, case, implication direction, or algebra?
\item Which three definitions must you review before starting Velleman?
\item Choose one proof from this exam and explain its architecture without symbols.
\item Choose one failed attempt and write a repair plan.
\item How will you structure your first two weeks with Velleman?
\item What will count as evidence that you are reading actively rather than passively?
\end{enumerate}

\subsection{Scoring guide}

Use this guide after your first attempt.

\begin{longtable}{p{0.18\linewidth}p{0.72\linewidth}}
\toprule
Score band & Interpretation \\
\midrule
85--100\% & Start Velleman now. Keep a proof-debugging log and continue spaced review. \\
70--84\% & Start Velleman slowly. Review weak topics every other day and rewrite missed exam proofs. \\
55--69\% & Delay full-speed Velleman. Spend two weeks on the readiness plan, especially definitions and proof skeletons. \\
Below 55\% & Return to the first-pass reader and rebuild the foundations before starting Velleman. \\
\bottomrule
\end{longtable}

\subsection{Final readiness contract}

Before you move on, write and sign a short contract with yourself:

\begin{quote}
I will not treat confusion as a verdict. I will treat it as a location. When a proof fails, I will find the first bad line, name the gap, open the relevant definition, and repair the argument.
\end{quote}

\begin{focusbox}
You are ready for the next book when you can work this way without needing someone else to rescue the proof at the first sign of fog.
\end{focusbox}
